\section{组合构建与交易摩擦}
\label{ch:lower-multifactor-research}

\begin{itemize}
  % Generated from curriculum/manifest.json; do not edit by hand.
\item 直接先修：下册《多因子估计：正则化、数值实现与外部数据》。

\end{itemize}

因子分数给出排序，交易需要确定数量。下面从同一组股票的排序出发，算出目标权重、实际换手和成本，让每一步都能追溯到前一步。


\MFQTransition{从预测到仓位与收益}

信号到净收益的联系可以写为
\[
x_{i,t}\longrightarrow \tilde x_{i,t}\longrightarrow q_{i,t}
\longrightarrow w_{i,t}\longrightarrow R^{\mathrm{gross}}_{t+1}
\longrightarrow R^{\mathrm{net}}_{t+1}.
\]
$\tilde x$ 是时点预处理后的信号，$q$ 是分组，$w$ 是实际仓位。组合收益由这些仓位与之后实现的资产收益共同决定；只计算信号的相关性，还不能得到组合收益。

% Generated by tools/build_figures.py; edit figures/figure-specs.json instead.
\MFQFigure{tex/figures/generated/lower-mf-pipeline.pdf}{本例依次将信号转为中性化暴露、目标权重与持仓；持仓产生收益，调仓带来成本。}{lower-mf-pipeline}


\MFQTransition{分组、权重与重叠持有期}

先在每个调仓日用事先确定资产池排序，再把最低组做空、最高组做多。等权使每一侧绝对权重和为 1；市值权重则在组内按正市值归一化。两者回答不同问题，不能因为某一结果更好而事后选择。

若持有 $H$ 期，每个时点的有效仓位是仍存续 vintage 的平均：
\[
w_t=\frac1{|A_t|}\sum_{s\in A_t}w_t^{(s)},
\qquad A_t=\{s:t-H+1\le s\le t\}.
\]
这会平滑仓位并引入收益重叠；推断时需处理序列相关。资产退市、停牌或不可卖空不能靠删除行解决，必须进入可交易集合和退出规则。

\MFQTransition{换手、成本与容量}

双边换手定义为 $\tau_t=\sum_i|w_{i,t}-w_{i,t-1}|$。净收益可写为
\[
R^{\mathrm{net}}_{t+1}=w_t^{\mathsf T}r_{t+1}
-c\tau_t-\operatorname{impact}(w_t,\mathrm{ADV}_t,\sigma_t).
\]
两期等权例的毛收益分别为 $0.035$ 与 $0.030$，期均毛收益 $0.0325$；平均双边换手为 2，单位换手成本 $0.001$，容量冲击 $0.0005$，所以净收益为 $0.0300$。市值加权与两期持有给出不同仓位和净收益，这是一项预先声明的敏感性分析，不是择优汇报。

% Generated by tools/build_figures.py; edit figures/figure-specs.json instead.
\MFQFigure{tex/figures/generated/lower-mf-cost-chain.pdf}{预测优势只有经过权重映射、换手和成本扣除后才成为净收益；每一层都可能改变排序。}{lower-mf-cost-chain}


真实报告应拆出佣金、价差、冲击、借券、税费、参与率、成交完成率与资金规模。容量不是一个永久常数；应报告规模到净收益、参与率和成交失败的曲线。

\MFQTransition{两种容易混淆的研究结论}

\textbf{情形 A：预测与收益断开。} 程序打印 AR、回归或 ML 预测，但 \texttt{trade\_return} 直接来自 合成样本。审计步骤是追踪每一笔 $w_t$ 是否由事先确定预测生成，并用 $w_t r_{t+1}$ 独立重算。修复不是给报告补一句“模型驱动”，而是删除外生收益字段。

\textbf{情形 B：选择后重写检验族。} 研究者测试 500 个信号，只把 20 个较好结果交给 BH，并把 $m$ 写成 20。检验族应包括实际参与选择的 500 个信号。只保留较好的 20 项会改变错误率计算；这与某项结果是否来自样本外是两个需要分别判断的问题。

\MFQTransition{分组组合的逐步算例}

假设五只股票得分排序为 $A>B>C>D>E$，未来收益为 $(3\%,1\%,0\%,-1\%,-2\%)$。最简单多空组合做多 $A$、做空 $E$，毛收益为 $5\%$；若每侧取两只等权，则收益为
\[
\frac{3\%+1\%}{2}-\frac{-1\%-2\%}{2}=3.5\%.
\]
单只头尾更容易被个股异常主导，分组更稳但信号稀释。市值加权、风险加权和分数加权回答不同投资问题，不能在测试期择优。

若下一期排名变成 $E>A>C>B>D$，目标仓位变化决定换手。换手应从旧实际仓位到新目标仓位计算，部分成交时还要保留未完成订单。只按入选名单变化计数会忽略权重大小。

\MFQTransition{例：从两次排名计算换手}

前一期排名为 $A>B>C>D>E$，每侧取两只等权，权重按 $A,B,C,D,E$ 排列为
\[
 (0.5,0.5,0,-0.5,-0.5).
\]
下一期排名变为 $E>A>C>B>D$，新权重为 $(0.5,-0.5,0,-0.5,0.5)$。暂忽略持有期间权重漂移，调仓差为 $(0,-1,0,0,1)$，双边换手为 2。

若单位换手成本为 10 个基点，成本为 $0.001\times2=0.002$。下一期收益为 $(0.03,0.01,0,-0.01,-0.02)$，新仓位毛收益为 $0.015-0.005+0.005-0.01=0.005$，净收益为 $0.003$。高分股票 $E$ 实现负收益，说明排序规则可以明确而单期预测仍失败。

真实调仓应先把上期持仓按已实现价格更新到调仓前权重，再计算订单；若只成交部分订单，下一期必须从实际仓位继续。notebook 将把两个排名转成权重数组并显示每只股票的贡献，读者可改变成本率，直接找出本期盈亏平衡成本 $0.005/2=0.0025$。

\MFQTransition{归因：模型收益究竟来自哪里}

组合收益可分为市场、行业、风格和特异部分。一个“质量 Alpha”若主要来自低波动或小盘暴露，应先说明它是复合暴露，而非把中性化后的剩余全部称为选股能力。归因至少按日期、行业、市值、流动性和多空两侧展开。

回撤归因比平均收益更能暴露失效：是信号翻转、拥挤、成本上升、股票池变化还是数据缺口？对每种失败预先定义诊断统计量，避免亏损后任意讲故事。

\MFQTransition{稳健性不是堆叠更多回测}

合理稳健性包括邻近参数、不同但事前合理的股票池、子时期、成本压力和数据供应商交叉核对。尝试越多，越要把全部结果计入选择权。热力图不能只截取正收益区域。

真正的最终测试是研究设计事先确定后一次性打开的未来区间。一旦利用这一结果继续修改模型，该区间就参与了选择；下一次确认需要新的、未参与修改的数据。

\begin{diagnostic}
高 IC、单调分组和显著 t 值可以同时存在，但若信号换手过高、容量过低或依赖不可借券股票，仍不是可实施策略。统计证据和实施证据必须并列。
\end{diagnostic}

\section{本路线的综合项目}

选择一个因子问题，将前面的回归、估计与组合构建联系起来。在未参与选择的区间比较分组收益，并考察成本与交易规模变化的影响。另选一种信号失效或成交受限的情形，解释净收益为何改变。

\section{综合练习}
\begin{enumerate}[leftmargin=*]
  \item \textbf{口述概念：}为什么预测型横截面系数、经典风险价格和因果效应不是同一对象？
  \item \textbf{口述概念：}何时 Pearson IC 与 Rank IC 会给出不同结论？
  \item \textbf{笔试推导：}从两期斜率 $0.020,0.030$ 推导均值和标准误。
  \item \textbf{笔试推导：}写出经典两遍法的两根回归轴，并解释生成回归量误差。
  \item \textbf{数值编程：}复现 beta 与风险价格，再把资产 alpha 改成不同值观察第二遍截距。
  \item \textbf{数值编程：}给零信号搜索加入相关候选，比较 BH 结论并记录依赖条件。
  \item \textbf{研究判断：}审计一份把预测系数称为风险价格的摘要，列出缺失假设。
  \item \textbf{研究判断：}说明为什么最终测试期不能参与信号尺度、窗口和检验族选择。
\end{enumerate}

\MFQLead{估计、正则化与数值实现题组}
\begin{enumerate}[leftmargin=*]
  \item \textbf{口述概念：}数值稳定、统计识别和经济解释分别由什么证据支持？
  \item \textbf{口述概念：}为什么“lasso 选中”不等于变量解释稳定？
  \item \textbf{笔试推导：}从 ridge 目标函数推导线性方程。
  \item \textbf{笔试推导：}推导 lasso 单坐标的软阈值更新。
  \item \textbf{数值编程：}画出 ridge/lasso 惩罚路径并比较透明实现与成熟库。
  \item \textbf{数值编程：}注入全零列与近共线列，分别记录失败或条件数响应。
  \item \textbf{研究判断：}审计一份用最终测试期选择 $\lambda$ 的研究包。
  \item \textbf{研究判断：}解释 WDI 截面为什么不能支持股票 alpha 或因果结论。
\end{enumerate}

\MFQTransition{组合实施、应用分析与开放项目 题组}
\begin{enumerate}[leftmargin=*]
  \item \textbf{口述概念：}信号、排序、目标权重、成交权重和净收益分别是什么对象？
  \item \textbf{口述概念：}等权、市值权重与重叠持有期改变了什么研究问题？
  \item \textbf{笔试推导：}手算两期等权组合的毛收益、换手、成本和净收益。
  \item \textbf{笔试推导：}写出 $H$ 期重叠持有的有效权重及其推断后果。
  \item \textbf{数值编程：}实现五分组、等权/市值权重与 1/3/6 期持有敏感性表。
  \item \textbf{数值编程：}从各期实际持仓和下一期资产收益计算组合毛收益，再根据成交增量扣除成本。
  \item \textbf{研究判断：}审计“只保留 20 个候选再做 BH”的研究报告。
  \item \textbf{研究判断：}审计未记录停牌、融券与成交完成率的 A 股组合。
  \item \textbf{综合项目：}选择一个因子，完成从回归到分组组合的分析。比较两种持有期或权重规则，并解释成本与样本选择怎样影响结果。
\end{enumerate}

\subsection*{基础衔接：独立计算}
同一组合的总权重与绝对权重和分别是多少？能否把 0.5\% 毛收益称作未加杠杆收益？

